\input{../tex-inputs/latex-leseansicht-vorspann}

\section[Gustav Schwarzkopf an Arthur Schnitzler, 20. 2. 1920]{L04258 Gustav Schwarzkopf an Arthur Schnitzler, 20. 2. 1920}
\nopagebreak\mylabel{L04258v}
\rehead{ }\normalsize\beginnumbering\briefempfaengerindex{Schnitzler, Arthur@\textsc{Schnitzler, Arthur}!zzzSchwarzkopf, Gustav@\emph{von Gustav Schwarzkopf}!1920-02-201@{20. 2. 1920}|(be}
\toendnotes[C]{\smallbreak\pagebreak[2]}
\correspDesc{Versand  durch Gustav Schwarzkopf am 20. 2. 1920 in Baden bei Wien
\newline{}Übermittlung  am 20. 2. 1920 in Baden bei Wien
\newline{}Erhalt  durch Arthur Schnitzler im Zeitraum [21. 2. 1920
                  – 25. 2. 1920?] in Wien}\toendnotes[C]{\smallbreak}
\Standort{CUL, Schnitzler, B 96a.}
\physDesc{Postkarte, 633 Zeichen
\newline{}Handschrift: schwarze Tinte, deutsche Kurrent
\newline{}Versand: Stempel: »\nobreak{}\oindex{Baden bei Wien@\textbf{Baden bei Wien}, \emph{Hauptstadt}|pwk}Baden b. Wien, 21 II 20, 5\nobreak{}«.  }\toendnotes[C]{\smallbreak}\pstart{}{\pb}\label{K_L04258-1v}\edtext{Z.}{\lemma{\textnormal{\emph{Z.}}}\Cendnote{\textnormal{Zimmer}}}\label{K_L04258-1} 15.\pend{}\pstart{}Städtiſche Curanſtalt\oindex{Städtische Kuranstalt@\textbf{Städtische Kuranstalt}, \emph{Beherbergungsgebäude}|pw}\pend{}{\bigskip}\pstart{}Herrn\pend{}\pstart{}\textsc{D\textsuperscript{r} Arthur Schnitzler}\pend{}\pstart{}\textsc{Wien (Cottage) XVIII\oindex{Wien@\textbf{Wien}!XVIII., Währing@\textbf{XVIII., Währing}!Währinger Cottage@\textbf{Währinger Cottage}, \emph{Teil eines besiedelten Ortes}|pw}}\pend{}\pstart{}Sternwarteſtraße 71\oindex{Wien@\textbf{Wien}!XVIII., Währing@\textbf{XVIII., Währing}!Sternwartestraße 71@\textbf{Sternwartestraße 71}, \emph{Wohngebäude}|pw}\pend{}{\bigskip}\vspace{1em}
\pstart
           \raggedleft{}{\pb}Baden\oindex{Baden bei Wien@\textbf{Baden bei Wien}, \emph{Hauptstadt}|pw}{ }20. II 20.\pend
           
\pstart{}Lieber Arthur,\pend\vspace{0.5em}
\pstart
           ich bin geſtern hier angeko{\geminationm}en und finde vorläufig alles{ }ſchön und vortrefflich. Auch das Wetter hatte einige Freundlichkeit für mich. Denken
               Sie nur, ich – \uline{ich} hab die So{\geminationn}e nach Baden\oindex{Baden bei Wien@\textbf{Baden bei Wien}, \emph{Hauptstadt}|pw}
               gebracht. Eine geeignetere Perſönlichkeit ſcheint ſie nicht gefunden zu haben. Baden\oindex{Baden bei Wien@\textbf{Baden bei Wien}, \emph{Hauptstadt}|pw} iſt noch i{\geminationm}er ſchön und ich hab auch Gelegenheit die Lücken in meiner Operettenbildung
               auszufüllen. Das Theater\oindex{Stadttheater Baden@\textbf{Stadttheater Baden}, \emph{Theater}|pwv} ſpielt täglich bei vollen Häuſern. – Wie geht es bei Ihnen? Nichts {\pb}Neues aus Berlin\oindex{Berlin@\textbf{Berlin}, \emph{Hauptstadt}|pw}?\pend
           
\pstart
           Ich grüße Frau \textsc{Olga\pwindex{Schnitzler, Olga 17.\,1.\,1882 Wien – 13.\,1.\,1970 Lugano@\textsc{Schnitzler, Olga} (17.\,1.\,1882 Wien – 13.\,1.\,1970 Lugano), \emph{Schauspielerin, Sängerin}|pw}},{\\[\baselineskip]}Sie und die Kinder\pwindex{Cappellini, Lili 13.\,9.\,1909 Wien – 26.\,7.\,1928 Venedig@\textsc{Cappellini, Lili} (13.\,9.\,1909 Wien – 26.\,7.\,1928 Venedig)|pwv}\pwindex{Schnitzler, Heinrich 9.\,8.\,1902 Hinterbrühl – 12.\,7.\,1982 Wien@\textsc{Schnitzler, Heinrich} (9.\,8.\,1902 Hinterbrühl – 12.\,7.\,1982 Wien), \emph{Regisseur, Schauspieler}|pwv}{\\[\baselineskip]}herzlichſt{\\[\baselineskip]}Ihr{\\[\baselineskip]}\spacefill\mbox{Gustav S}\pend
           \leftskip=0em{}\selectlanguage{ngerman}\endnumbering\briefempfaengerindex{Schnitzler, Arthur@\textsc{Schnitzler, Arthur}!zzzSchwarzkopf, Gustav@\emph{von Gustav Schwarzkopf}!1920-02-201@{20. 2. 1920}|)be}\mylabel{L04258h}
\begin{anhang}
\end{anhang}\newcommand{\dateiname}{L04258}\newcommand{\titel}{Gustav Schwarzkopf an Arthur Schnitzler, 20. 2. 1920}\newcommand{\editorInnen}{Herausgegeben von Jahnke, SelmaMüller, Martin Anton}\input{../tex-inputs/latex-leseansicht-abspann}
